\documentclass[11pt]{article}
\usepackage{amsmath,amsfonts}
\usepackage{graphicx}
%\usepackage{xcolor}
%\definecolor{gray}{rgb}{.75,.75,.75}
%\usepackage[landscape]{geometry}
%\usepackage{hyperref}
%\usepackage[bookmarks=true, pdfborder={0 0 .5}, urlbordercolor=gray]{hyperref}

\renewcommand{\thefootnote}{\fnsymbol{footnote}}


\addtolength{\topmargin}{-.6in}
\addtolength{\textheight}{1.2in}
\addtolength{\oddsidemargin}{-.3in}
\addtolength{\textwidth}{.6in}
\pagestyle{empty}

\newcommand{\ds}{\displaystyle}
\newcommand{\ts}{\textstyle}

\newcommand{\Z}{\mathbb{Z}}
\newcommand{\R}{\mathbb{R}}
\newcommand{\C}{\mathbb{C}}
\newcommand{\EE}{\mathcal{E}}
\newcommand{\tZ}{\tilde{\Z}}

\newcommand{\Sym}{\mathrm{Sym}}

\renewcommand{\circle}[1]{{\bigcirc\hspace{-.75em}#1}}

\begin{document}

\footnote[0]{Notes by Peter Magyar \ {\tt magyar@math.msu.edu}}
\vspace{-1em}

\centerline{\bf Math 133\hspace{1.5em} \hfill{\Large\bf
Comparison Tests}\hfill Stewart \S11.4}
\vspace{1em}

\noindent 
{\bf Convergence and divergence.} We continue to discuss convergence tests: ways to tell if a given series $\sum_{n=1}^\infty a_n=\lim_{N\to\infty} \sum_{n=1}^N a_n$ converges (to a finite value),
or diverges (to infinity or by oscillating).\footnote{
A general divergent series might oscillate up and down forever, but a positive series (with $a_n\geq 0$) either levels off to a finite value, or diverges to infinity.
}
So far, we know convergence for two kinds of standard series:
\begin{itemize}

\item Geometric series: $\sum_{n=1}^\infty cr^{n-1}$ converges to $\frac{c}{1-r}$ if $|r|<1$, diverges if $|r|\geq 1$.

\item Standard $p$-series: $\sum_{n=1}^\infty \frac{1}{n^p}$ converges if $p>1$, and diverges if $p\leq 1$.

\end{itemize}
%
%We also know two general tests: the $n$th Term Vanishing Test and the Integral Test.
In this section, we test convergence of a complicated series $\sum a_n$ by comparing it to a simpler one 
(such as the above): a convergent ceiling $\sum c_n$, or a divergent floor~$\sum d_n$.
\\[1em]
{\bf  Direct Comparison Test:} Let $M$ be a positive integer starting point.
\begin{itemize}

\item If $0\,{\leq}\, a_n\,{\leq}\, c_n$ for $n\,{\geq}\,M$, and $\sum\limits_{n=1}^\infty c_n$ converges, then $\sum\limits_{n=1}^\infty a_n$ converges.

\item If $a_n\,{\geq}\, d_n\,{\geq}\, 0$ for $n\,{\geq}\,M$, and $\sum\limits_{n=1}^\infty d_n$ diverges, then $\sum\limits_{n=1}^\infty a_n$ diverges.
\end{itemize}
%
These results are clear, since the series $\sum_{n=1}^\infty a_n$ is term-by-term smaller or larger than its comparison series, except possibly the first $M{-}1$ terms.\footnote{
Here we use the completeness axiom of real analysis: if a series of increasing partial sums 
has an upper bound, $s_N=\sum_{n=1}^N a_n<B$ for all $N$, then the limit (least upper bound)
$L=\lim\limits_{N\to \infty} s_N$ exists.
}
\\[.5em]
{\sc Example:} Determine convergence of: \  
$ \ds\sum\limits_{n=1}^\infty a_n =\sum\limits_{n=1}^\infty \frac{n-1}{n^2\sqrt n \,+\,1}$.
\\
We compare to the series $c_n$ obtained from the dominant terms of the top and bottom,
and use the general principles for inequality of fractions:
$
\frac{\text{smaller top}}{\text{larger bot}} \leq \frac{\text{top}}{\text{bot}} \leq \frac{\text{larger top}}{\text{smaller bot}}.
$
$$
a_n=\frac{n-1}{n^2\sqrt n\,+\,1}
\ \ \leq\ \ c_n=\frac{n}{n^2\sqrt n}=\frac1{n^{3/2}}
\qquad\text{for $n\geq 1$,}
$$
since in $c_n$ the top is larger and the bottom smaller. 
The ceiling series $\sum_{n=1}^\infty c_n = \sum_{n=1}^\infty \frac1{n^{3/2}}$ 
is a standard $p$-series which converges, so $\sum_{n=1}^\infty a_n$ also converges.
\\[.5em]
{\sc Example:} Determine convergence of: \  
$\ds\sum\limits_{n=1}^\infty  \frac{2^{3n+\sin(n)}}{3^n + 4n^2}$.
\\
As a rough guess, we ignore the lower-order terms in numerator and denominator to compare with $\frac{2^{3n}}{3^n}
\ =\ \left(\frac{8}{3}\right)^{\!n}$, which makes a divergent geometric series, so our series $a_n$ should also diverge.
However, $d_n$ is not of the form $\frac{\text{smaller top}}{\text{larger bot}}$, so we cannot use 
$\sum_{n=1}^\infty d_n=\sum_{n=1}^\infty\left(\frac{8}{3}\right)^{\!n} = \infty$ 
as a divergent floor for $a_n$ in the the Comparison Test.

We need a smaller $d_n$, with smaller top and larger bottom. 
To decrease the top: $2^{3n+\sin(n)}=
2^{3n}2^{\sin(n)}\geq 2^{3n}2^{-1}$.
To increase the bottom, we take a slightly larger geometric series:
in fact
$4^n\geq 3^n+4n^2$ for all $n\geq 3$.
Thus:
$$
a_n = \frac{2^{3n+\sin(n)}}{3^n + n^2}
\ \ \geq\ \ 
d_n = \frac{2^{3n}2^{-1}}{4^n}
=\tfrac12 2^n 
\qquad\text{for $n\geq 3$.}
$$
We only need this inequality for all large $n$: 
the first couple of terms $a_1, a_2$ make no difference 
to convergence or divergence of the infinite sums.
Since $\sum_{n=1}^\infty d_n = \sum_{n=1}^\infty\tfrac12 2^n=\infty$ is a divergent floor, 
the orginal $\sum_{n=1}^\infty a_n$ also diverges.
\\[.5em]
{\sc Example:} Determine convergence of: \  
$\ds\sum\limits_{n=1}^\infty \frac{n+1}{n^3 - 20}$\,.
\\
Roughly estimate top and bottom by their leading terms:
$\sum_{n=1}^\infty\frac{n}{n^3}=\sum_{n=1}^\infty\frac{1}{n^2}$, a convergent standard $p$-series,
so the original series probably converges.
But $a_n=\frac{n+1}{n^3-20}>\frac n{n^3}$, so we cannot use $c_n=\frac n{n^3}$ as a convergent ceiling for $a_n$.
But maybe:
$$
a_n=\frac{n+1}{n^3-20}\ \leq\ c_n=2\frac{n}{n^3} \qquad
\text{for $n$ large enough.}
$$
How large does $n$ need to be to make this inequality valid?
Let us check:
$$
\frac{n+1}{n^3-20}\leq\frac2{n^2}
\quad\Longleftarrow\quad
0<n^2(n+1)\leq 2(n^3-20)
\quad\Longleftrightarrow\quad
40\leq n^2(n{-}1)
\quad\Longleftarrow\quad
n\geq 4\,.
$$
Thus, we have:
$$
a_n=\frac{n+1}{n^3-20}\ \leq\ c_n = \frac2{n^2} \qquad
\text{for $n\geq 4$,}
$$
where $\sum_{n=1}^\infty\frac{2}{n^2} = 2\sum_{n=1}^\infty\frac{1}{n^2}$ converges,
so the original $\sum_{n=1}^\infty a_n$ also converges.
\\[.5em]
{\sc example:} Consider any infinite decimal:
$$
s \ =\ 0.d_1d_2d_3\cdots \ =\ \frac{d_1}{10}+\frac{d_2}{10^2}+\frac{d_3}{10^3}+\cdots
\ =\ \sum\limits_{n=1}^\infty \frac{d_n}{10^n}\,,
$$
where $0\leq d_n\leq 9$ are {\it any} decimal digits. Does this series always converge,
so that the infinite decimal represents a real number, or could a bad choice of digits define a meaningless decimal? 

In fact, we can compare $0\leq \frac{d_n}{10^n}\leq \frac{9}{10^n}$, 
since each digit is at most 9.
The ceiling is a convergent geometric series: $\sum_{n=1}^\infty \frac{9}{10^n} =\sum_{n=1}^\infty \frac9{10}\left(\frac1{10}\right)^{\!n-1} = \frac{9}{10}\frac{1}{1-\frac1{10}}=1$, so the original decimal sequence also converges.
{\it Any} infinite decimal represents a number.

\pagebreak

\noindent
{\bf Limit Comparison Test.}
Suppose $\lim_{n\to\infty} \frac{a_n}{b_n} = L$ with $0<L<\infty$.
\begin{itemize}

\item If $\sum\limits_{n=1}^\infty b_n$ converges, then $\sum\limits_{n=1}^\infty a_n$ converges.

\item If $\sum\limits_{n=1}^\infty b_n$ diverges, then $\sum\limits_{n=1}^\infty a_n$ diverges.
\end{itemize}
%
{\small 
{\it Proof:} $\lim\limits_{n\to\infty} \frac{a_n}{b_n} = L$ means that, 
for any small $\epsilon>0$, 
we can take a starting point $N$ so that  for all $n\geq N$, we have:
$$
L{-}\epsilon \ \leq\ \tfrac{a_n}{b_n} \ \leq\ L{+}\epsilon
\qquad\text{and}\qquad
(L{-}\epsilon) b_n \ \leq\ a_n \ \leq\ (L{+}\epsilon) b_n\,.
$$ 
Taking $\epsilon$ small enough that $L{\pm}\epsilon>0$, 
we can prove convergence or divergence
by taking $c_n=(L{+}\epsilon) b_n$ or
$d_n=(L{-}\epsilon) b_n$ in the Direct Comparison Test.
}
\\[.5em]
{\sc Example:} We redo \  
$\ds\sum\limits_{n=1}^\infty \frac{n+1}{n^3 - 20}$\,.
Now we can immediately compare with $b_n =\dfrac{n}{n^3}$\,:
$$
\frac{a_n}{b_n}
\ =\
\frac{n+1}{n^3 - 20}\,\left/\frac{n}{n^3}\right.
\ =\
\frac{n+1}{n}\left/\frac{n^3-20}{n^3}\right.
\ =\
\frac{1+\frac1n}{1-\frac{20}{n^3}}\,.
$$
Taking $n\to\infty$ gives $L=1$. 
Since this satisfies $0<L<\infty$,  
and $\sum_{n=1}^\infty b_n=\sum_{n=1}^\infty \frac1{n^2}$ 
is a convergent standard $p$-series, the original series
$\sum_{n=1}^\infty a_n$ also converges.
\\[1em]
{\bf Extended Limit Comparison Test.} Now let $a_n,b_n> 0$ be positive series.
In the case where $\lim\limits_{n\to\infty}\frac{a_n}{b_n}=L=0$,
we have $a_n$ much smaller than $b_n$, so if $\sum_{n=1}^\infty b_n<\infty$ converges, then so does  $\sum_{n=1}^\infty a_n<\sum_{n=1}^\infty b_n<\infty$. 
Similarly, in the case where $\lim\limits_{n\to\infty}\frac{a_n}{b_n}=L=\infty$,
we have $a_n$ much larger than $b_n$, so if $\sum_{n=1}^\infty b_n=\infty$ diverges, then so does  $\sum_{n=1}^\infty a_n>\sum_{n=1}^\infty b_n=\infty$.
\\[.5em]
{\sc example:} Determine the convergence of:  $\ds\sum_{n=1}^\infty \frac{n^2}{2^n}$\,. 
\\
Since $n^2$ is negligeable compared to the exponential growth of $2^n$,
we could roughly estimate this by $\sum_{n=1}^\infty b_n = \sum_{n=1}^\infty \frac{1}{2^n}
= \sum_{n=1}^\infty \left(\frac 12\right)^{\!n}$, 
a convergent geometric series, so the original series should converge.

However, taking the Limit Comparison Test with this $b_n=\frac{1}{2^n}$ gives 
$L=\infty$, since $a_n=\frac{n^2}{2^n}$ is much {\it larger} than $b_n$.
Thus this comparison fails: 
$b_n$ is a convergent floor for $a_n$, and we can't tell whether the larger $\sum a_n$ converges or diverges.

We take a slightly larger, but still convergent, comparison: $b_n=\left(\frac 34\right)^{\!n}$:
$$
\lim_{n\to\infty} \frac{a_n}{b_n} 
\ =\ 
\lim_{n\to\infty}  \frac{ n^2\!\left(\frac12\right)^{\!n} } {\left(\frac 34\right)^{\!n} } 
\ =\ 
\lim_{n\to\infty} n^2 \!\left( \tfrac23 \right)^{\!n}
\ =\ 0\,,
$$
as we could prove by L'H\^opital's Rule. Thus $a_n=\frac{n^2}{2^n}$ becomes much {\it smaller}
than $b_n$, and $\sum_{n=1}^\infty b_n=\sum_{n=1}^\infty \left(\frac 34\right)^{\!n}<\infty$ is a
convergent ceiling for $\sum_{n=1}^\infty a_n$, which therefore must also converge.
\\[.5em]

\end{document}

%%%%%%%%%%%  Picture  %%%%%%%%%%
\\[1em]
\centerline{
%\includegraphics[width=2in]
{1-8-Discont-iv.png}
\qquad \qquad 
%\includegraphics[width=2in]
{1-8-Discont-v.png} 
}
\vspace{0em}

\noindent
%%%%%%%%%%%%%%%%%%%%%%%%%

%%%%%%%%%%%  Table  %%%%%%%%%%
In fact, we get the following derivatives:
$$\begin{array}{|c||c|c|c|c|c|c|}
\hline &&&&&& \\[-.9em] 
f(x) & \sin(x) & \cos(x) & \tan(x)  & \sec(x) & \csc(x) & \cot(x)
\\[.2em] \hline &&&&&& \\[-.9em]
f'(x) &\cos(x) & -\sin(x) & \sec^2(x) & \tan(x)\sec(x) & -\cot(x)\csc(x) & -\csc^2(x)
\\[.1em] \hline
\end{array}$$
\\[1em]
%%%%%%%%%%%%%%%%%%%%%%%%%%



















